% !PRINTED 501
\weberpage{501}
%<<CONT
value, then it is an integral function of $j(\omega)$. If in (5)
we let $x$ pass over into $v$, there results
\weq{6}{M = \frac{\psi(v,u)}{F'(v)},}
in which $\psi(v,u)$ is an integral rational function of $u$ and
$v$, with rational numerical coefficients.

If we now let $\omega$ in (6) pass over into the root of a
quadratic equation whose discriminant
\[
D = B^2 - 4AC
\]
is negative:
\weq{7}{A\omega^2 + B\omega + C = 0,}
then one or more of the values $v$ can become equal to $u$. If
this is the case for only one of the $\nu$ values $v$, then $F'(v)$
will not vanish for $v = u$, and by (6) we can express $M$
rationally through this singular value of $u$.

But the equation $u = v$ carries with it the condition:
\[
\frac{c+\partial\omega}{a} = \frac{\gamma+\delta\omega}{\alpha+\beta\omega}.
\]
From this it follows:
\weq{8}{\begin{aligned}
\partial\beta &= Ax, & \partial\alpha &= \frac{Bx-y}{2},\\
c\beta-a\delta &= \frac{Bx+y}{2}, & c\alpha-a\gamma &= Cx,
\end{aligned}}
\weq{9}{4n = y^2 - Dx^2,}
in which $x$ can be assumed positive.

We assume $A$ odd and relatively prime to $D$. If
$D \equiv 1\ (\mathrm{mod}\ 4)$, we put $n = -D$ and obtain from (9):
\[
y^2 + n(x^2-4) = 0.
\]
According to this, the value $x=1$ is admissible only when $-D$
is three times a square number. If we set this case aside for
the moment, there remains only
\[
x = 2, \quad y = 0,
\]
and from (8) there results:
\[
\begin{aligned}
\partial &= 1, & \beta &= 2A, & \alpha &= B,\\
a &= n, & \gamma &\equiv 2C, & \delta &\equiv B\ (\mathrm{mod}\ 4),\\
\alpha + \beta\omega &= \sqrt{D}, &&&&
\end{aligned}
\]
and by the last two equations (8) $c$ is determined to the
modulus $n$.
